{\actuality} 
% компоненты оптимизатора
Повышение эффективности выполнения запросов остаётся одной из центральных задач теории и практики систем управления базами данных (СУБД). Скорость выполнения запроса в базе данных в значительной степени определяется тем, насколько удачно оптимизатор выбирает план выполнения из множества допустимых альтернатив. На этот выбор влияют оценки числа строк в промежуточных результатах запроса, то есть оценки кардинальностей, стоимостная модель, которая приближённо оценивает затраты альтернативных планов, и алгоритм перебора пространства планов. Неточность хотя бы в одном из этих компонентов может приводить к выбору существенно менее эффективного плана и, как следствие, к заметному росту времени выполнения запроса \cite{10.14778/2850583.2850594}.

% тестируемся на аналитичечкой нагрузке
Особенно остро данная проблема проявляется при аналитических нагрузках. Для такого типа нагрузки характерны сложные многотабличные запросы, включающие большое число операций соединения и фильтрации, а также обработку значительных объёмов данных. Такие запросы обычно формируют широкое пространство альтернативных планов выполнения, в котором итоговая эффективность выбранного плана существенно зависит от точности оценок кардинальностей, адекватности стоимостной модели и эффективности алгоритма перебора пространства планов. Кроме того, при аналитических нагрузках влияние неточностей оптимизатора выражено особенно сильно, поскольку даже локально неудачное решение на одном из этапов плана может приводить к заметному росту общего времени выполнения запроса \cite{10.14778/2850583.2850594, 10.1145/3318464.3380584}.

Таким образом, актуальной является разработка методов повышения качества планирования запросов, которые сохраняют существующий оптимизатор СУБД, но позволяют учитывать фактические условия выполнения и направлять выбор альтернативных планов при приемлемых накладных расходах.

{\progress}
% кардинальность - развитое направление



Одним из наиболее подробно исследованных направлений в области интеллектуальной оптимизации запросов являются методы повышения точности оценок кардинальностей \cite{10.14778/3485450.3485459, 10.1145/3514221.3526154}. Для этой задачи предложен широкий спектр подходов~--- от классических статистических методов \cite{10.1145/376284.375686} до современных обучаемых моделей \cite{kipf2018learnedcardinalitiesestimatingcorrelated}. В то же время даже существенное повышение точности оценок кардинальностей не устраняет полностью проблему выбора неоптимального плана, поскольку на итоговое качество оптимизации влияют также стоимостная модель и ограничения алгоритма перебора пространства планов \cite{10.14778/2850583.2850594, 10.14778/1687627.1687739}. Вопросы оценки кардинальностей исследуются в работах N.~Bruno, S.~Chaudhuri, L.~Gravano, A.~Kipf, J.~Sun, K.~Kim и других авторов \cite{10.1145/376284.375686,kipf2018learnedcardinalitiesestimatingcorrelated,10.14778/3485450.3485459,10.1145/3514221.3526154}.

% интеграция с существующим оптимизатором
При этом полная замена штатного оптимизатора внешней моделью в практических СУБД связана с высокой инженерной сложностью и не всегда целесообразна. Современный оптимизатор содержит большое число правил преобразования запросов, эвристик, механизмов учёта физических операторов, индексов, статистики и параметров среды выполнения. Его замена требует воспроизведения значительной части этой функциональности, а также обеспечения корректности, устойчивости и переносимости решений при изменении данных, схемы, рабочей нагрузки и конфигурации системы. Поэтому практически значимым направлением являются гибридные методы, которые сохраняют существующий оптимизатор, но улучшают отдельные компоненты процесса планирования или направляют выбор альтернативных планов через предусмотренные СУБД механизмы управления \cite{zhu2023lero, anneser2023autosteer}.

% стоимостная модель



Одной из таких точек интеграции является стоимостная модель оптимизатора. Существующие исследования показывают, что повышение точности стоимостной модели может улучшать качество выбора плана оптимизатором. На практике это достигается как за счёт калибровки параметров стоимостной модели, так и за счёт обучаемых моделей оценки планов, интегрируемых в оптимизатор \cite{Lan_2021}. Однако такие решения обычно требуют накопления обучающих данных, регулярного переобучения при изменении нагрузки и нередко плохо переносятся между различными конфигурациями систем. Кроме того, ручная настройка параметров стоимости требует высокой квалификации администратора базы данных, а калибровка на специально сформированной синтетической нагрузке не всегда отражает реальные условия эксплуатации. Поэтому актуальной остаётся задача адаптивной настройки параметров стоимостной модели по телеметрии фактически выполняемых запросов, без запуска отдельной калибровочной нагрузки, без обучения отдельной модели \cite{10.1145/3514221.3526154, 10.1145/3318464.3380584}. Вопросам построения и настройки стоимостных моделей посвящены работы T.~Siddiqui и других исследователей; состояние данного направления обобщено H.~Lan, Z.~Bao и Y.~Peng \cite{10.1145/3318464.3380584,Lan_2021}.

% перебор пространства планов
Наряду с оценками кардинальностей и стоимостной моделью существенное влияние на итоговое качество оптимизации оказывает алгоритм перебора пространства планов. Из-за высокой вычислительной сложности полного перебора оптимизаторы используют различные ограничения и эвристики, позволяющие удерживать время планирования в разумных пределах, но способные приводить к пропуску части потенциально эффективных альтернатив \cite{Lan_2021}. Поэтому отдельный практический интерес представляют методы, позволяющие направлять выбор планов без полной замены существующего оптимизатора.

% подсказки оптимизатору



Одним из практических средств направленного воздействия на поведение оптимизатора являются подсказки, позволяющие задавать дополнительные ограничения или предпочтения при построении плана выполнения и тем самым влиять на множество рассматриваемых альтернатив и итоговый план выполнения. В последние годы активно развиваются подходы, основанные на машинном обучении, в которых механизм выбора подсказок используется для управления существующим оптимизатором без его полной замены. Вместе с тем автоматический подбор подсказок остаётся сложной задачей: пространство возможных конфигураций может быть достаточно большим, полный перебор требует значительных затрат, ошибочная подсказка способна ухудшить план, а обучаемые решения часто требуют предварительного накопления данных и дополнительного сопровождения \cite{10.1145/3448016.3452838, zhu2023lero}. Методы обучаемого выбора планов и управления штатным оптимизатором посредством подсказок развиваются, в частности, в работах R.~Marcus, R.~Zhu, C.~Anneser и их соавторов \cite{10.1145/3448016.3452838,zhu2023lero,anneser2023autosteer}.

% режимы автоматического подбора подсказок
В задачах автоматического подбора подсказок можно выделить два практически различных режима, различающихся допустимыми накладными расходами и целью оптимизации. Первый режим связан с онлайн-оптимизацией, когда рекомендация должна быть получена непосредственно перед выполнением запроса, а время выбора подсказки должно быть существенно меньше ожидаемого выигрыша от её применения. Такой сценарий характерен для потоков аналитических запросов, где требуется быстро принять решение для нового запроса на основе уже накопленного опыта, не выполняя дорогостоящий перебор и не запуская дополнительное обучение на целевой нагрузке. 

Второй режим связан с офлайн-оптимизацией или предварительной настройкой, когда для запроса или группы похожих запросов допустимо выделить ограниченный бюджет пробных запусков. Такой сценарий оправдан, например, для часто повторяющихся аналитических запросов, ускорение которых будет использоваться многократно. В этом случае основная цель состоит не в мгновенной рекомендации, а в нахождении более эффективной конфигурации подсказок оптимизатору. Поэтому онлайн- и офлайн-постановки соответствуют разным практическим задачам: первая ориентирована на минимальные накладные расходы при принятии решения, вторая~--- на более тщательный поиск качественного решения при заранее заданном бюджете оптимизации.

% языковые модели для БД
%Дополнительный интерес связан с быстрым развитием больших языковых моделей. В настоящее время большие языковые модели в области баз данных применяются по нескольким крупным направлениям. Во-первых, они используются для построения интерфейсов на естественном языке к данным, включая преобразование запросов на естественном языке в SQL и поддержку диалогового взаимодействия с базой данных \cite{Zhou2024DBGPTLL,10.1007/s11704-025-50592-w}. Во-вторых, они применяются для обработки структурированных представлений базы данных, включая анализ таблиц, схем таблиц и связанных с ними метаданных, а также для сопоставления схем и интеграции источников данных \cite{10.1007/s11704-024-40763-6,parciak2024schemamatchinglargelanguage}. В-третьих, языковые модели используются для задач сопровождения и настройки СУБД, включая интеллектуальные средства работы с документацией \cite{Zhou2024DBGPTLL,10.1007/s00778-023-00831-y}. В-четвёртых, активно исследуется применение больших языковых моделей для оптимизации запросов, включая переписывание запросов, а в отдельных работах - и непосредственное построение планов выполнения \cite{li2024llmr2,song2026quitequeryrewriterules,10.1145/3626246.3654732}. Вместе с тем существующие подходы на основе больших языковых моделей имеют ряд ограничений. К их числу относятся зависимость качества результата от выбора подсказок и примеров, слабый учёт физических характеристик СУБД и данных, трудности обеспечения эквивалентности преобразований и повышенные риски галлюцинаций. Дополнительно в обзорах отмечаются сохраняющиеся проблемы обработки сложных запросов, устойчивости, безопасности и вычислительной стоимости таких систем \cite{Zhou2024DBGPTLL, song2026quitequeryrewriterules, 10.1007/s11704-025-50592-w}.

% языковые модели для БД



Дополнительный интерес к указанным задачам связан с быстрым развитием больших языковых моделей и расширением области их применения в СУБД. В настоящее время такие модели исследуются для построения интерфейсов на естественном языке к данным, обработки таблиц, схем и метаданных, сопровождения и настройки СУБД, а также для задач оптимизации запросов, включая переписывание запросов и построение вариантов планов выполнения \cite{Zhou2024DBGPTLL,10.1007/s11704-025-50592-w,10.1007/s11704-024-40763-6,parciak2024schemamatchinglargelanguage,10.1007/s00778-023-00831-y,li2024llmr2,song2026quitequeryrewriterules,10.1145/3626246.3654732}. Это показывает, что большие языковые модели становятся одним из перспективных инструментов для задач, связанных с анализом, настройкой и автоматизацией работы СУБД. Поэтому актуальным является исследование способов их использования не только в пользовательских интерфейсах или генерации SQL-запросов, но и в задачах повышения качества планирования запросов. Применение больших языковых моделей к задачам сопровождения СУБД, переписывания запросов и построения планов рассматривается в работах I.~Trummer, Z.~Li, Y.~Song, M.~Urban, C.~Binnig и других авторов \cite{10.1007/s00778-023-00831-y,li2024llmr2,song2026quitequeryrewriterules,10.1145/3626246.3654732}.

% режимы использования языковых моделей
При этом способ применения больших языковых моделей должен соответствовать ограничениям конкретной постановки оптимизации. В режиме онлайн-рекомендации основным ограничением является время принятия решения: подсказку необходимо выбрать до выполнения запроса, а накладные расходы на оптимизацию должны быть малы. Поэтому в такой постановке целесообразно использовать языковую модель в облегчённом режиме, извлекая из неё компактное представление, пригодное для быстрого сопоставления нового запроса с ранее накопленным опытом. В режиме офлайн-оптимизации или предварительной настройки, напротив, допустим ограниченный бюджет пробных запусков, а основной целью является получение наиболее эффективной найденной конфигурации подсказок оптимизатору. В этом случае языковая модель может использоваться более полно, как компонент, принимающий решения в итеративном процессе оптимизации.

% актуальность исследования
%Таким образом, анализ степени разработанности темы показывает, что к настоящему времени наиболее подробно исследованы методы оценки кардинальностей и отдельные подходы, основанные на обучении, для выбора планов и подсказок. Вместе с тем в меньшей степени разработаны подходы, связанные с онлайн-настройкой стоимостной модели по телеметрии реальных запросов без обучения отдельной модели. Существенный интерес также представляют быстрые методы рекомендации подсказок, использующие  большие языковые модели, но не требующие длительной генерации и специального обучения. При этом применение больших языковых моделей в задачах оптимизации и сопровождения баз данных всё ещё связано с рядом открытых вопросов. К ним относятся учёт физических характеристик конкретной СУБД и статистики данных, проверка корректности предлагаемых решений, снижение задержек инференса и стоимости дообучения, а также переносимость методов между различными схемами данных, рабочими нагрузками и конфигурациями вычислительной среды.

% актуальность исследования

Таким образом, анализ степени разработанности темы показывает, что к настоящему времени наиболее подробно исследованы методы оценки кардинальностей и отдельные обучаемые подходы к оценке и выбору планов. Вместе с тем остаются открытыми вопросы настройки стоимостной модели по телеметрии реальных запросов без отдельной обучаемой модели. Существенный интерес также представляют методы рекомендации подсказок, использующие возможности больших языковых моделей без необходимости длительной генерации при обработке каждого запроса. Применение больших языковых моделей в задачах оптимизации и сопровождения баз данных всё ещё связано с рядом открытых вопросов, включая учёт физических характеристик конкретной СУБД и статистики данных, проверку корректности предлагаемых решений, снижение задержек инференса и стоимости дообучения, а также переносимость методов между различными схемами данных, рабочими нагрузками и конфигурациями вычислительной среды.


{\object} Планирование запросов в реляционных системах управления базами данных.

{\subject} Методы повышения качества планирования запросов в реляционных системах управления базами данных.

{\aim} данной работы является разработка и экспериментальное обоснование комплекса методов повышения качества планирования запросов в реляционных системах управления базами данных, включающего адаптивную настройку параметров стоимостной модели, рекомендацию подсказок оптимизатору на основе векторных представлений планов выполнения, формируемых большой языковой моделью, и агентный итеративный поиск подсказок.

Для~достижения поставленной цели необходимо решить следующие {\tasks}:
\begin{enumerate}[beginpenalty=10000] % https://tex.stackexchange.com/a/476052/104425
  \item Провести анализ современного состояния методов оптимизации запросов в реляционных СУБД и определить ограничения существующих подходов, связанные с настройкой стоимостной модели, управлением выбором альтернативных планов и применением обучаемых методов в составе существующего оптимизатора.
  \item Разработать метод адаптивной настройки параметров стоимостной модели, обеспечивающий адаптацию к фактическим условиям выполнения запросов по телеметрии эксплуатационной нагрузки, без запуска отдельной калибровочной нагрузки и без обучения отдельной модели. Оценить влияние метода на согласованность между оценочной стоимостью и фактическим временем выполнения.
  \item Разработать метод рекомендации подсказок оптимизатору для онлайн-сценария, в котором выбор управляющего воздействия должен выполняться с малыми накладными расходами на основе ранее накопленного опыта, без обучения на целевой рабочей нагрузке. Исследовать возможность использования векторных представлений, формируемых большой языковой моделью, для сопоставления планов выполнения и использования ранее найденных подсказок для новых запросов.
  \item  Разработать метод итеративного поиска подсказок оптимизатору для офлайн-сценария или предварительной настройки, в котором допустим ограниченный бюджет пробных запусков, а выбор очередного действия выполняется с учётом текущего состояния поиска. Исследовать эффективность применения большой языковой модели в качестве компонента принятия решений в таком контуре.
  \item Провести экспериментальную оценку разработанных методов на репрезентативных аналитических наборах запросов и сравнить их с соответствующими базовыми и альтернативными подходами по времени выполнения, качеству выбора планов и стоимости оптимизации.
\end{enumerate}

{\novelty}
\begin{enumerate}[beginpenalty=10000] % https://tex.stackexchange.com/a/476052/104425
\item Разработан метод адаптивной настройки параметров стоимостной модели оптимизатора, отличающийся использованием агрегированной модели состояния буферного кэша на уровне отношений и настройкой параметров по телеметрии фактически выполняемых запросов. Метод позволяет уточнять параметры ввода-вывода и вычислительной составляющей стоимости без отдельной калибровочной нагрузки и без замены штатного оптимизатора.

\item Разработан метод рекомендации подсказок оптимизатору на основе сопоставления планов выполнения в пространстве векторных представлений. В методе предобученная языковая модель используется для формирования векторных представлений планов выполнения, а выбор кандидата выполняется на основе сопоставления этих представлений, после чего выполняется дополнительная проверка согласованности. Метод позволяет использовать ранее накопленный опыт выполнения запросов без дообучения модели на целевой рабочей нагрузке.

\item Разработан агентный метод итеративного поиска подсказок оптимизатору в замкнутом контуре обратной связи с системой управления базами данных. Метод  основан на использовании большой языковой модели для пошагового выбора действий и повышает результативность поиска эффективных конфигураций подсказок при ограниченном бюджете пробных запусков.
\end{enumerate}

{\methods} В диссертационной работе для решения поставленных задач использовались методы теории баз данных, методы машинного обучения, методы математической статистики, методы математического и компьютерного моделирования, методы регрессионного анализа, методы анализа данных, а также методы вычислительного эксперимента. Экспериментальные исследования выполнялись на аналитических наборах запросов с применением сравнительного анализа времени выполнения, корреляционных характеристик стоимостной модели и показателей качества выбора планов выполнения.

{\theoreticalSignificance} Предложен и обоснован комплекс методов повышения качества планирования SQL-запросов в реляционных СУБД на основе адаптации стоимостной модели, сопоставления планов выполнения в пространстве векторных представлений и итеративного поиска подсказок оптимизатору. Данный подход развивает методы гибридной оптимизации запросов, в которых средства машинного обучения дополняют штатный оптимизатор без его полной замены.

{\influence} заключается в том, что предложенные методы могут быть использованы для повышения качества выбора планов выполнения в СУБД в рамках существующего оптимизатора. Разработанные решения ориентированы на применение в аналитических нагрузках, допускают внедрение с ограниченными изменениями в системе и позволяют уменьшить объём ручной настройки параметров и перебора конфигураций подсказок. Практическая значимость исследования подтверждается актом о внедрении результатов работы.

{\defpositions}
\begin{enumerate}[beginpenalty=10000] % https://tex.stackexchange.com/a/476052/104425
\item Метод адаптивной настройки параметров стоимостной модели оптимизатора по телеметрии фактически выполняемых запросов, учитывающий состояние буферного кэша, позволяет повышать согласованность между стоимостью плана и фактическим временем его выполнения, а также улучшать качество выбора плана без отдельного этапа калибровки.

\item Метод рекомендации подсказок оптимизатору на основе сопоставления планов выполнения, использующий векторные представления планов, формируемые большой языковой моделью, позволяет использовать ранее накопленный опыт для выбора эффективных конфигураций подсказок без обучения на целевой рабочей нагрузке.

\item Агентный метод итеративного поиска подсказок оптимизатору на основе большой языковой модели позволяет направленно искать эффективные конфигурации подсказок в замкнутом контуре обратной связи с системой управления базами данных и повышать качество оптимизации при ограниченном бюджете пробных запусков по сравнению с базовыми и альтернативными методами.
\end{enumerate}

{\pasport} Тема диссертационного исследования и его содержание соответствуют требованиям паспорта специальности ВАК 2.3.5 – Математическое и программное обеспечение вычислительных систем, комплексов и компьютерных сетей:
\begin{enumerate}
    \setcounter{enumi}{2}
    \item Модели, методы, архитектуры, алгоритмы, языки и программные инструменты организации взаимодействия программ и программных систем.
    \item Интеллектуальные системы машинного обучения, управления базами данных и знаний, инструментальные средства разработки цифровых продуктов. 
\end{enumerate}

{\reliability} полученных результатов обеспечивается использованием стандартных и широко применяемых в исследовательской практике бенчмарков для оценки методов оптимизации запросов, сопоставлением предложенных методов с базовыми и альтернативными подходами, проведением экспериментов в контролируемых условиях, а также использованием воспроизводимой методики оценки, включающей повторные измерения и анализ совокупных показателей времени выполнения и качества принимаемых решений.

{\probation}
Основные результаты работы докладывались на:
\begin{enumerate}
    \item семинар Московской секции ACM SIGMOD в 2026 году, семинар лаборатории информационных систем ИСИ СО РАН в 2025 году, семинар русскоязычного сообщества AGI в 2025 году;
    \item Международная конференция <<Знания -- Онтологии -- Теории>> (ЗОНТ), 2025;
    \item IEEE XVII International Conference on Actual Problems of Electronic Instrument Engineering (APEIE), 2025.
\end{enumerate}
%Основные результаты работы докладывались на:
%перечисление основных конференций, симпозиумов и~т.\:п.

{\contribution} Работа выполнена в составе научной группы. Исходные идеи и постановки задач, рассматриваемых в главах 2–4 диссертации, развивались в неразрывном сотрудничестве с соавторами. Автором лично выполнены программная реализация предложенных методов, подготовка данных, планирование и проведение вычислительных экспериментов, обработка и анализ полученных результатов, а также подготовка текста диссертации. В главе 2 автором самостоятельно разработана агрегированная модель состояния буферного кэша, выполнены программная реализация метода адаптивной настройки параметров стоимостной модели, проведение экспериментальных исследований и анализ результатов; формализация части метода, связанной с настройкой параметров процессора, выполнена совместно с научным руководителем. В главе 3 автором самостоятельно разработаны предсказательная модель и метод рекомендации подсказок оптимизатору на основе сопоставления планов выполнения, выполнены их реализация, подготовка экспериментальных данных, проведение экспериментов и анализ результатов. В главе 4 разработка агентного метода итеративного поиска подсказок оптимизатору, формирование контекста и диалогового набора данных, реализация агентов, проведение экспериментов и анализ результатов выполнены автором самостоятельно.

\ifnumequal{\value{bibliosel}}{1}
{%%% Встроенная реализация с загрузкой файла через движок bibtex8. (При желании, внутри можно использовать обычные ссылки, наподобие `\cite{vakbib1,vakbib2}`).
    \nocite{vasilenko2024acm}
    \nocite{vasilenko2026plan_mapping}
    \nocite{vasilenko_article_inpress}
    \nocite{11289347}
    \nocite{vasilenko2025}
    \nocite{burl2025}
    \nocite{vasilenko_patent}
    {\publications} Основные результаты диссертационного исследования отражены в 6 научных публикациях, включая 2 статьи в журналах, индексируемых в международной базе данных Scopus, из которых 1 статья опубликована в журнале первого квартиля и 1 статья~--- в журнале второго квартиля по SCImago Journal Rank; 1 статью в журнале, рекомендованном ВАК; а также 3 работы, опубликованные в сборниках и трудах международных и всероссийских научных конференций, из которых 1 работа индексируется в базе данных Scopus. Получены свидетельство о государственной регистрации программы для ЭВМ и акт о внедрении результатов исследования.
}%
{%%% Реализация пакетом biblatex через движок biber
    \begin{refsection}[bl-author, bl-registered]
        % Это refsection=1.
        % Процитированные здесь работы:
        %  * подсчитываются, для автоматического составления фразы "Основные результаты ..."
        %  * попадают в авторскую библиографию, при usefootcite==0 и стиле `\insertbiblioauthor` или `\insertbiblioauthorgrouped`
        %  * нумеруются там в зависимости от порядка команд `\printbibliography` в этом разделе.
        %  * при использовании `\insertbiblioauthorgrouped`, порядок команд `\printbibliography` в нём должен быть тем же (см. biblio/biblatex.tex)
        %
        % Невидимый библиографический список для подсчёта количества публикаций:
        \phantom{\printbibliography[heading=nobibheading, section=1, env=countauthorvak,          keyword=biblioauthorvak]%
        \printbibliography[heading=nobibheading, section=1, env=countauthorwos,          keyword=biblioauthorwos]%
        \printbibliography[heading=nobibheading, section=1, env=countauthorscopus,       keyword=biblioauthorscopus]%
        \printbibliography[heading=nobibheading, section=1, env=countauthorconf,         keyword=biblioauthorconf]%
        \printbibliography[heading=nobibheading, section=1, env=countauthorother,        keyword=biblioauthorother]%
        \printbibliography[heading=nobibheading, section=1, env=countregistered,         keyword=biblioregistered]%
        \printbibliography[heading=nobibheading, section=1, env=countauthorpatent,       keyword=biblioauthorpatent]%
        \printbibliography[heading=nobibheading, section=1, env=countauthorprogram,      keyword=biblioauthorprogram]%
        \printbibliography[heading=nobibheading, section=1, env=countauthor,             keyword=biblioauthor]%
        \printbibliography[heading=nobibheading, section=1, env=countauthorvakscopuswos, filter=vakscopuswos]%
        \printbibliography[heading=nobibheading, section=1, env=countauthorscopuswos,    filter=scopuswos]}%
        %
        \nocite{*}%
        %
        {\publications} Основные результаты по теме диссертации изложены в~\arabic{citeauthor}~печатных изданиях,
        \arabic{citeauthorvak} из которых изданы в журналах, рекомендованных ВАК%
        \ifnum \value{citeauthorscopuswos}>0%
            , \arabic{citeauthorscopuswos} "--- в~периодических научных журналах, индексируемых Web of~Science и Scopus%
        \fi%
        \ifnum \value{citeauthorconf}>0%
            , \arabic{citeauthorconf} "--- в~тезисах докладов.
        \else%
            .
        \fi%
        \ifnum \value{citeregistered}=1%
            \ifnum \value{citeauthorpatent}=1%
                Зарегистрирован \arabic{citeauthorpatent} патент.
            \fi%
            \ifnum \value{citeauthorprogram}=1%
                Зарегистрирована \arabic{citeauthorprogram} программа для ЭВМ.
            \fi%
        \fi%
        \ifnum \value{citeregistered}>1%
            Зарегистрированы\ %
            \ifnum \value{citeauthorpatent}>0%
            \formbytotal{citeauthorpatent}{патент}{}{а}{}%
            \ifnum \value{citeauthorprogram}=0 . \else \ и~\fi%
            \fi%
            \ifnum \value{citeauthorprogram}>0%
            \formbytotal{citeauthorprogram}{программ}{а}{ы}{} для ЭВМ.
            \fi%
        \fi%
        % К публикациям, в которых излагаются основные научные результаты диссертации на соискание учёной
        % степени, в рецензируемых изданиях приравниваются патенты на изобретения, патенты (свидетельства) на
        % полезную модель, патенты на промышленный образец, патенты на селекционные достижения, свидетельства
        % на программу для электронных вычислительных машин, базу данных, топологию интегральных микросхем,
        % зарегистрированные в установленном порядке.(в ред. Постановления Правительства РФ от 21.04.2016 N 335)
    \end{refsection}%
    \begin{refsection}[bl-author, bl-registered]
        % Это refsection=2.
        % Процитированные здесь работы:
        %  * попадают в авторскую библиографию, при usefootcite==0 и стиле `\insertbiblioauthorimportant`.
        %  * ни на что не влияют в противном случае

        %\nocite{progbib1}%program
        %\nocite{bib1}%other
        %\nocite{confbib1}%conf
    \end{refsection}%
        %
        % Всё, что вне этих двух refsection, это refsection=0,
        %  * для диссертации - это нормальные ссылки, попадающие в обычную библиографию
        %  * для автореферата:
        %     * при usefootcite==0, ссылка корректно сработает только для источника из `external.bib`. Для своих работ --- напечатает "[0]" (и даже Warning не вылезет).
        %     * при usefootcite==1, ссылка сработает нормально. В авторской библиографии будут только процитированные в refsection=0 работы.
}


